This problem is a continuation of the linear regression problem from the previous homework (HW1 Problem 5). Let us recall the setup. Given
$d$-dimensional input data
$\textbf{x}_1,\ldots, \textbf{x}_n \in \mathbb{R}^d$ with
real-valued labels $y_1,\ldots, y_n \in \mathbb{R}$, the goal
is to find the coefficient vector $\textbf{w}$ that minimizes the sum
of the squared errors.  The total squared error of $\mathbf{w}$ can be
written as
$f(\mathbf{w})=\sum_{i=1}^{n}f_i(\mathbf{w})$, where
$f_i(\mathbf{w})=(\mathbf{w^T x_i}-y_i)^2$ denotes the
squared error of the $i$th data point. We will refer to $f(\w)$ as the \emph{objective function} for the problem.

Last homework, we considered the scenario where the number of data points was much larger than the number of dimensions and
hence we did not worry too much about generalization. (If you remember, the gap between training and test accuracies in 5.1 of the last HW was not too large).  We will
now consider the setting where $d=n$, and examine the test error along
with the training error. Use the following Python code for generating
the training data and test data.

\begin{verbatim}
import numpy as np

train_n = 100
test_n = 1000
d = 100

X_train = np.random.normal(0,1, size=(train_n,d))
w_true = np.random.normal(0,1, size=(d,1))
y_train = X_train.dot(w_true) + np.random.normal(0,0.5,size=(train_n,1))

X_test = np.random.normal(0,1, size=(test_n,d))
y_test = X_test.dot(w_true) + np.random.normal(0,0.5,size=(test_n,1))

\end{verbatim}

\textbf{3.1} (\blue{2pts}) We will first setup a baseline, by finding the test
  error of the linear regression solution
  $\textbf{w}=\mathbf{X^{-1}y}$ without any regularization.  This is
  the closed-form solution for the minimizer of the objective function
  $f(\mathbf{w})$.  (Note the formula is simpler than what we saw in the last homework because
  now $\mathbf{X}$ is square as $d=n$).  Report the training error and test
  error of this approach, \emph{averaged over 10 trials}. For better
  interpretability, report the normalized error
  $\hat{f}(\textbf{w})$ rather than the value of the objective function
  ${f}(\textbf{w})$, where we define $\hat{f}(\textbf{w})$ as
\[ 
\hat{f}(\textbf{w})=\frac{\parallel{\mathbf{Xw-y}}\parallel_2}{\parallel{\mathbf{y}}\parallel_2}.
\]

\emph{Note on averaging over multiple trials:} We're doing this to get a better estimate of the performance of the algorithm. To do this, simply run the entire process (including data generation, and sampling $w_{\text{true}}$) 10 times, and find the average value of $\hat{f}(\textbf{w})$ over these 10 trials. \\

\noindent{\fontsize{11}{12}\selectfont\textbf{Solution 3.1}

For the unregularized least-squares objective
\[
F(\mathbf w)=\|X\mathbf w-\mathbf y\|_2^2,
\]
setting has \(d=n\), so we use the closed form
\[
\mathbf w_{LS}=X^{-1}\mathbf y.
\]
For each trial, we compute \(\hat f\) on train and test data, then average across 10 trials. The resulting values are:
\[
\hat f_{train}=3.6466\times 10^{-15},
\qquad
\hat f_{test}=2.9623.
\]
The training error is essentially zero while the test error is large. This indicates interpolation of training noise and clear overfitting in this high-dimensional regime. The matrix is an almost necessarily invertible square matrix, so the closed-form solution actually precisely solves the linear system, this indicates that the model can perfectly fit the training data (including the noise within it), so the training error is numerically close to the machine precision error. Although the training error is almost zero, due to the same number of parameters and samples, the model not only fits the actual signal but also captures the noise in the training data. Therefore, its generalization ability is very poor on new test data.

\bigskip


\textbf{3.2} (\blue{7pts}) We will now examine $\ell_2$ regularization as a means to prevent overfitting. The $\ell_2$ regularized objective function is given by the following expression:
$$ \sum_{i=1}^{m} (\mathbf{w^T x_{i}}-{y}_{i})^2 + \lambda \| \mathbf{w}\|_2^2.$$
As discussed in class, this has a closed-form solution
$\textbf{w}=\mathbf{(X^TX+\lambda I)^{-1}X^Ty}$. Using this closed-form
solution, present a plot of the normalized training error and
normalized test error $\hat{f}(\textbf{w})$ for
$\lambda = \{0.0005, 0.005, 0.05, 0.5, 5, 50, 500\}$. As before, you
should average over 10 trials. Discuss the characteristics of your
plot, and also compare it to your answer to~(\textbf{3.1}).\\

\noindent \emph{The following questions explore the concept of implicit regularization. This is a very active topic of research, with the idea being that optimization algorithms such as SGD can themselves act like regularizers (in the sense that they prefer the solutions to the regularized problems instead of just the original problem). There's no one correct answer we're looking for in many of these questions, and idea is to make you think about what is happening and report your findings.}\\

\noindent{\fontsize{11}{12}\selectfont\textbf{Solution 3.2}

We now solve ridge regression:
\[
\min_{\mathbf w}\ \|X\mathbf w-\mathbf y\|_2^2+\lambda\|\mathbf w\|_2^2,
\qquad
\mathbf w_{ridge}(\lambda)=(X^TX+\lambda I)^{-1}X^T\mathbf y.
\]
We evaluate
\[
\lambda\in\{0.0005,0.005,0.05,0.5,5,50,500\}
\]
and report averages over 10 trials.

\begin{table}[H]
\centering
\caption{Averaged normalized train/test errors for ridge regression}
\begin{tabular}{ccc}
\toprule
$\lambda$ & Avg. train error & Avg. test error \\
\midrule
0.0005 & 0.0016 & 0.7017 \\
0.005  & 0.0029 & 0.4241 \\
0.05   & 0.0060 & 0.2663 \\
0.5    & 0.0155 & \textbf{0.2324} \\
5      & 0.0651 & 0.3416 \\
50     & 0.2784 & 0.5691 \\
500    & 0.7300 & 0.8567 \\
\bottomrule
\end{tabular}
\end{table}

The best test performance is obtained at \(\lambda=0.5\), with test error \(0.2324\).

\begin{figure}[H]
\centering
\includegraphics[width=0.78\linewidth]{fig_3_2_ridge_curve.png}
\caption{Train and test error versus $\lambda$}
\end{figure}

As \(\lambda\) increases, the training error monotonically rises, while the test error shows a typical U-shaped curve. When \(\lambda\) is small, the model almost perfectly fits the training data (the training error is close to 0), but the test error is relatively high; when \(\lambda\) is moderate (about 0.5), the test error reaches its minimum; as \(\lambda\) further increases, both the training and test errors rise simultaneously, indicating underfitting. 

When \(\lambda\) is very small, regularization has almost no effect and the model approaches a non-regularized state, thus still overfitting the training noise and resulting in a high test error. An appropriate \(\lambda\) will penalize the weight norm, making the model smoother and thereby improving the generalization performance; however, when \(\lambda\) is too large, the model is overly constrained and cannot adequately fit the true signal, leading to underfitting and the test error rising again.

Comparison with 3.1 
In 3.1, the test error without regularization was approximately 6.09, which was much higher than the test error of about 0.22 when using the optimal \(\lambda\) (approximately 0.5) in this problem. This indicates that $\ell_2$ regularization significantly alleviated the overfitting problem and greatly enhanced the generalization ability of the model. In contrast, without regularization, although the training error was close to 0, the test error was extremely high, demonstrating severe overfitting.

\bigskip

\textbf{3.3} (\blue{7pts}) Run stochastic gradient descent (SGD) on the original objective
function $f(\mathbf{w})$, with the initial guess of $\textbf{w}$ set to be the all 0's vector.  Run SGD for 1,000,000 iterations for each different choice of the step size, $\{0.00005$, $0.0005$, $0.005\}$.  Report the
 normalized training error and the normalized test
error for each of these three settings, averaged over 10 trials. How does the SGD solution compare with the
solutions obtained using $\ell_2$ regularization? Note that SGD is
minimizing the original objective function, which does \emph{not} have any
regularization. In Part (\textbf{3.1}) of this problem, we found the
\emph{optimal} solution to the original objective function with
respect to the training data. How does the training and test error of
the SGD solutions compare with those of the solution in~(\textbf{3.1})? 
Can you explain your observations?   (It may be helpful to also compute the normalized training and test error corresponding to the true coefficient vector $\textbf{w}^*$ for comparison, this is \verb|w_true| in the code.)\\

\noindent{\fontsize{11}{12}\selectfont\textbf{Solution 3.3}

We run SGD on the original (unregularized) objective. For one sampled point \((\mathbf x_i,y_i)\), let
\[
F_i(\mathbf w)=(\mathbf w^T\mathbf x_i-y_i)^2,
\quad
\nabla F_i(\mathbf w)=2(\mathbf w^T\mathbf x_i-y_i)\mathbf x_i.
\]
Hence the update is
\[
\mathbf w^{(t+1)}=\mathbf w^{(t)}-\eta\cdot 2(\mathbf w^{(t)T}\mathbf x_i-y_i)\mathbf x_i.
\]
We use \(\mathbf w^{(0)}=\mathbf 0\), run for \(10^6\) iterations, and average over 10 trials for
\(\eta\in\{0.00005,0.0005,0.005\}\).

\begin{table}[H]
\centering
\caption{Averaged normalized train/test errors for SGD}
\begin{tabular}{ccc}
\toprule
Step size $\eta$ & Avg. train error & Avg. test error \\
\midrule
0.00005 & 0.0135 & 0.2148 \\
0.0005  & 0.0068 & \textbf{0.2041} \\
0.005   & 0.0034 & 0.4039 \\
\bottomrule
\end{tabular}
\end{table}

\begin{figure}[H]
\centering
\includegraphics[width=0.78\linewidth]{fig_3_3_norm_curve.png}
\caption{Training error trajectory, with horizontal reference line for $\hat f(w^*)$}
\end{figure}

\begin{figure}[H]
\centering
\includegraphics[width=0.78\linewidth]{fig_3_3_train_test_curve.png}
\caption{Training error trajectory, with horizontal reference line for $\hat f(w^*)$}
\end{figure}

The best final test error is obtained at \(\eta=0.0005\), with value \(0.2041\).
As the step size increases, both the training error and the test error decrease significantly. A smaller step size ($5e^{-5}$) did not converge adequately after 1,000,000 iterations, resulting in high training and test errors; while a larger step size ($5e^{-3}$) could approach the optimal solution more quickly, thereby significantly reducing the errors. 
\bigskip

Compared with the optimal ridge solution in Section 3.2 (with a test error of approximately 0.22), the best test error of SGD under the current settings (approximately 0.40) is still higher. This indicates that although SGD performs well in optimizing the original objective function, due to the lack of explicit regularization, its generalization performance is still inferior to the solution with appropriate $\ell_2$ regularization. 

In comparison in 3.1, the training error of the closed-form solution is close to 0, but the test error is as high as approximately 6, indicating severe overfitting. However, the training error of the current SGD does not approach 0, and the test error is significantly lower. This suggests that SGD has not fully converged to the exact interpolation solution of the training data but has remained near a solution with a smaller norm, thereby demonstrating better generalization ability.

\bigskip

The optimization algorithm itself plays a role similar to regularization to some extent, and thus the obtained solutions perform better on the test set than the closed-form solution that perfectly fits the training data.

\bigskip
\textbf{3.4} (\blue{10pts}) We will now examine the behavior of SGD in more detail. Let $\mathbf{w}^{(t)}$ refer to the SGD solution at the $t$-th iteration. For step sizes $\{0.00005,0.005\}$ and 1,000,000 iterations of SGD (and a single trial), 
\begin{itemize}
\item [(i)] Plot the normalized training error of the SGD iterate $\mathbf{w}^{(t)}$ vs.\ the iteration
  number $t$. On the plot of training error, draw a line parallel to the
  x-axis indicating the error $\hat{f}(\textbf{w}^*)$ of
  the true model $\textbf{w}^*.$

	\item [(ii)] Plot the normalized test error of the SGD iterate $\mathbf{w}^{(t)}$ vs.\ the iteration number $t$. Your code might take a long time to run if you compute the test error after every SGD step---feel free to compute the test error every 100 iterations of SGD to make the plots.
	\item [(iii)] Plot the $\ell_2$ norm of the SGD iterate $\mathbf{w}^{(t)}$ vs.\
 the iteration number $t$.

\end{itemize}
  Comment on the plots. What can you say about the generalization ability of SGD with different step sizes? Does the plot correspond to the intuition that a learning algorithm starts to overfit when the training error becomes too small, i.e. smaller than the noise level of the true model so that the model is fitting the noise in the data? How does the generalization ability of the final solution depend on the $\ell_2$ norm of the final solution?\\

  \noindent{\fontsize{11}{12}\selectfont\textbf{Solution 3.4}
  
For a single trial, we run SGD for \(10^6\) iterations with step sizes \(\{0.00005,0.005\}\), and track trajectories (test metrics logged every 100 iterations).

\begin{figure}[H]
\centering
\includegraphics[width=0.78\linewidth]{fig_3_4_train_curve.png}
\caption{Training error trajectory, with horizontal reference line for $\hat f(w^*)$}
\end{figure}

The training error decreases monotonically with the increase in the number of iterations and eventually approaches 0, indicating that SGD can eventually fit the training data. However, with a larger step size, the training error decreases more rapidly, while with a smaller step size, it takes a longer time to reach the same level.

\begin{figure}[H]
\centering
\includegraphics[width=0.78\linewidth]{fig_3_4_test_curve.png}
\caption{Test error trajectory}
\end{figure}

The test error initially decreases as the training progresses, but after the training error continues to decrease, the test error may start to rise or stabilize. This phenomenon indicates that when the model begins to overfit the noise in the training data, its generalization ability will decline. Larger step sizes usually achieve a lower test error in the early stage, but they may also cause greater fluctuations.

\begin{figure}[H]
\centering
\includegraphics[width=0.78\linewidth]{fig_3_4_norm_curve.png}
\caption{Parameter norm trajectory $\|w^{(t)}\|_2$}
\end{figure}

As the iterations proceed, the weight norm $\|w\|_2$ continues to increase and reaches a maximum value when the training error approaches 0. This indicates that the model gradually increases the parameter norm in an attempt to further reduce the training error, which may lead to fitting noise. Smaller norms typically correspond to better generalization performance.

\begin{figure}[H]
\centering
\includegraphics[width=0.78\linewidth]{fig_3_4_norm_error_iter.png}
\caption{different steps sgd curves}
\end{figure}

As can be seen in the figure, when the training error becomes extremely small (close to 0), the test error no longer decreases but starts to increase, and at the same time, $\|w\|_2$ significantly increases. This is in line with the intuition of overfitting: when the training error is lower than the actual noise level, the model begins to fit the noise. When the test error is at its minimum, $\|w\|_2$ is usually small. As $\|w\|_2$ increases, the test error rises. This is consistent with the ridge result in 3.2: the small-norm solution has better generalization ability. Therefore, due to its limited number of steps and initialization method, SGD tends to converge to solutions with smaller norms, which is a form of implicit regularization.

\bigskip

\textbf{3.5} (\blue{5pts}) We will now examine the effect of the starting point on the SGD solution. Fixing the step size at $0.00005$ and the maximum number of iterations at 1,000,000, choose the initial point randomly from the $d$-dimensional sphere with radius $r=\{0, 0.1, 0.5, 1, 10, 20, 30\}$ (to do this random initialization, you can sample from the standard Gaussian $N(0,\mathbf{I})$, and then renormalize the sampled point to have $\ell_2$ norm $r$). Plot the average normalized training error and the average normalized test error over 10 trials vs $r$. Comment on the results, in relation to the results from part (\textbf{3.2}) where you explored different $\ell_2$ regularization coefficients. Can you provide an explanation for the behavior seen in this plot?\\

\noindent{\fontsize{11}{12}\selectfont\textbf{Solution 3.5}
We fix \(\eta=0.00005\) and \(10^6\) iterations, and initialize from a random point on the sphere of radius
\[
r\in\{0,0.1,0.5,1,10,20,30\}.
\]
All values are averaged over 10 trials.

\begin{table}[H]
\centering
\caption{Averaged normalized train/test errors versus initialization radius}
\begin{tabular}{ccc}
\toprule
Radius $r$ & Avg. train error & Avg. test error \\
\midrule
0   & 0.0116 & 0.2355 \\
0.1 & 0.0116 & 0.2353 \\
0.5 & 0.0116 & 0.2345 \\
1   & 0.0115 & \textbf{0.2338} \\
10  & 0.0132 & 0.2778 \\
20  & 0.0189 & 0.4064 \\
30  & 0.0260 & 0.5650 \\
\bottomrule
\end{tabular}
\end{table}

\begin{figure}[H]
\centering
\includegraphics[width=0.78\textwidth]{fig_3_5_init_radius_curve.png}
\caption{Train/test error versus initialization radius}
\end{figure}

\begin{figure}[H]
\centering
\includegraphics[width=0.78\textwidth]{fig_3_5_error_radius_curve.png}
\caption{Train/test error versus step size}
\end{figure}

As the initialization radius r increases, the training error slightly rises, while the test error significantly increases. When r is very small (close to 0), the test error is the lowest. When r increases to 10, 20, and 30, the test error significantly deteriorates. This indicates that the initialization scale has an impact on the generalization performance of the final model.

The reasons for these phenomena are that although SGD optimizes the unregularized objective function, the initialization radius actually affects the norm of the final converged solution. When starting from a smaller radius, SGD is more likely to converge to a solution with a smaller norm; while starting from a larger radius, the weights are more likely to maintain a larger norm during training, which leads to more severe overfitting.

\bigskip

Compared with 3.2, in 3.5, the smaller initialization radius plays a similar role - it limits the scale of the final solution of the model. Therefore, the initialization radius and $\ell_2$ act to some extent as an "implicit regularization" mechanism.

comparing with 3.4, it can be observed that different initialization radii can optimize the path, thereby affecting the final weight norm and consequently influencing the generalization ability.

\bigskip

3.5 It may also demonstrate an important fact: The convergence point of SGD is not only dependent on the objective function, but also on the optimization path and the initialization. When the initialization radius is very large, SGD needs to "shrink back" to a reasonable area; if the step size is small or the number of iterations is limited, the model may remain near a solution with a larger norm, thereby exhibiting poor test error.

\bigskip
\section*{Conclusion}
Across all subproblems, the same pattern appears: exact training fit does not guarantee good test performance. Explicit regularization (ridge) and implicit regularization (through SGD step size, trajectory, and initialization) both control effective complexity and improve generalization.


\bigskip
